\documentclass[11pt,letterpaper]{article}

\usepackage[T1]{fontenc}
\usepackage[utf8]{inputenc}
\usepackage[margin=1in]{geometry}
\usepackage{booktabs}
\usepackage{longtable}
\usepackage{array}
\usepackage{amsmath}
\usepackage{enumitem}
\usepackage{etoolbox}
\usepackage[htt]{hyphenat}
% Appendix A's provenance table has long results/*.log file-path cells in narrow p{} columns --
% allow \texttt to break at reasonable points instead of overflowing (matches the mitigation
% used for paper/pandoc-header.tex's earlier markdown-to-PDF build).
\AtBeginEnvironment{longtable}{\footnotesize\sloppy}
\emergencystretch=3em
\makeatletter
\let\origttfamily\ttfamily
\renewcommand{\ttfamily}{\origttfamily\hyphenchar\font=`\-\relax}
\makeatother
\usepackage[colorlinks=true,linkcolor=black,citecolor=black,urlcolor=blue,
  pdftitle={Judge-Shaped, Not Harm-Shaped: A Persistent Measurement-Validity Failure in
    Jailbreak Success Judges, and an Honest Null Result for Activation-Guided Mutation},
  pdfauthor={Muhammed Muiz Arummal}]{hyperref}
\usepackage[capitalize,noabbrev]{cleveref}
\usepackage[natbib=true,style=authoryear,backend=biber,maxcitenames=2,mincitenames=1]{biblatex}
\addbibresource{references.bib}

\crefname{table}{Table}{Tables}
\Crefname{table}{Table}{Tables}
\crefname{section}{Section}{Sections}
\Crefname{section}{Section}{Sections}
\crefname{subsection}{Section}{Sections}
\Crefname{subsection}{Section}{Sections}

\title{Judge-Shaped, Not Harm-Shaped: A Persistent Measurement-Validity Failure in Jailbreak
  Success Judges, and an Honest Null Result for Activation-Guided Mutation}
\author{Muhammed Muiz Arummal \\ \small Independent Researcher}
\date{DRAFT --- pre-arXiv --- reviewed three times (\texttt{reviews/stage7.md}), Gate 7 signed
  off 2026-08-08, numerically re-verified and converted to LaTeX 2026-08-09 (this file compiles
  clean: zero LaTeX errors, zero unresolved \texttt{\textbackslash cref}/\texttt{\textbackslash
  cite} warnings)\thanks{Not yet posted or tagged; disclosure to MSRC/Alibaba
  (\cref{sec:ethics}) is planned for posting time, not yet sent. See
  \texttt{reviews/stage7-human-signoff.md} for the sign-off record.}}

\begin{document}
\maketitle

\begin{abstract}
Jailbreak fuzzing benchmarks are usually reported as a single number: attack success rate
(ASR), scored by an automated judge. That number is only as trustworthy as the judge producing
it. We report a measurement-validity failure in \texttt{hubert233/GPTFuzz}, GPTFuzzer's default
RoBERTa judge. It rewards jailbreak-\textbf{shaped} surface form---persona/roleplay
declarations, ``you are now X'' framing---rather than harmful content. Hand-reading its flagged
``successes'' on Phi-4-mini-instruct found refusals that pivoted to an unrelated topic, and
completions that merely echoed a DAN/Omega-style template's setup with zero harmful content.

We built a stricter two-stage replacement: a deterministic refusal pre-filter plus an
anti-roleplay LLM-as-judge rubric. \textbf{The fix did not fully work.} On a fresh run, the
corrected judge still passed a ``ChadGPT'' persona-wrapper completion that actually refuses and
offers crisis-support resources---zero harmful content. Judge unreliability on persona-wrapper
completions survives a deliberate, targeted fix; it is not one classifier's quirk.

As a secondary result, we report an honest null. Using the corrected judge, activation-guided
span mutation---mutating only the template span with the highest projection onto a
difference-in-means refusal direction, rather than the whole template---shows no consistent,
reliable ASR advantage over GPTFuzzer's uniform-mutation baseline, on either target we tested
(Qwen2.5-3B-Instruct, Phi-4-mini-instruct). We independently confirmed the guided-mutation
mechanism was engaged: attribution fired on every iteration, and search revisited mutated
candidates. Attribution \textit{quality}, however, was template- and target-dependent, not
uniformly refusal-localized---consistent with a guidance signal that only sometimes helps.

All results are smoke-scale ($n=5$ behaviors per condition); the planned 25-behavior
$\times$ 3-seed matrix never ran---compute (Kaggle's free tier) was exhausted first. Both
findings argue that jailbreak-fuzzing evaluation infrastructure, not just attack methods,
deserves scrutiny.
\end{abstract}

% The provenance note below is part of the contribution (Section 1, contribution 3), not
% front-matter boilerplate. A standard-format venue may want it moved to supplementary material.
\begin{quote}
\noindent\textbf{How to read this paper.} Every number below is either (a) read directly from a
committed \texttt{results/*.json} / \texttt{results/*.npz} file---cited inline as
\texttt{path} $\to$ \texttt{field}---or (b) explicitly marked in-text as hand-verified (a
completed human check not written back into any file) or as narrative-only (no backing artifact
exists in this repository). Claims of the second kind are named as such wherever they appear
rather than being softened into unmarked prose, and no conclusion rests on one.
\end{quote}

\section{Introduction}
\label{sec:intro}

Attack success rate (ASR)---the fraction of harmful behaviors for which a fuzzing loop finds at
least one prompt an LLM complies with---is the standard headline metric in jailbreak-fuzzing
work, following GPTFuzzer \citep{yu2023gptfuzzer}. ASR is only as meaningful as the judge that
decides ``did this completion comply.'' That judge is frequently a small, frozen, off-the-shelf
classifier reused across many follow-on papers without re-validation---an implicit assumption
that the judge itself is a solved problem.

This paper started as a study of a different question: does mutating a jailbreak template at the
\textit{location} an interpretability signal points to (a refusal direction extracted via
difference-in-means over paired harmful/harmless prompts) produce better guided mutations than
GPTFuzzer's uniform, whole-prompt mutation? In the course of running that experiment, the
project's own pipeline caught its judge lying to it. Four completions the judge scored $\approx
0.99$ (``jailbroken'') on a Phi-4-mini smoke run were, on inspection, not jailbreaks at all---one
was a refusal that pivoted to an unrelated topic, three were the target model merely reciting a
roleplay template's setup with no harmful content whatsoever. The judge was detecting
jailbreak-\textbf{shaped} vocabulary, not harm.

We treat this not as an incident to patch quietly and move past, but as the paper's primary
finding. We built a stricter two-stage judge (\cref{sec:method}) and validated that it changes
real outcomes (\cref{sec:judge-reliability}). But hand-verifying the \textit{fixed} judge's own
output on a later, methodologically cleaner run turned up a second, distinct instance of the
same failure class---a persona-wrapper completion that actually refuses and provides crisis
resources, still scored a pass. That a deliberate, rubric-guided LLM-as-judge fix did not fully
close this gap is, we argue, more informative than either finding alone: judge reliability on
persona-wrapper/roleplay-shaped completions is a persistent problem across judge architectures,
not a quirk of one RoBERTa classifier.

Our secondary contribution is an honest null result, but a negative finding only earns trust if
its mechanism actually ran. \Cref{sec:mechanism} confirms guided mutation's attribution and
tree-search components were both genuinely engaged throughout every reported run---not silently
reduced to uniform mutation under the hood. Only against that confirmed-active mechanism does the
null become informative: using the corrected judge, guided mutation still shows no consistent,
reliable ASR advantage over uniform mutation on either target (\cref{sec:null}---we ran no formal
significance test at this sample size; ``no advantage'' here means the raw counts do not
consistently favor either method, not a statistically established equivalence). We can therefore
rule out ``the mechanism never engaged'' as the explanation, distinguishing this null from an
artifact of guided mutation quietly falling back to uniform behavior.

Both results come with an honest scope limitation: everything here is smoke-scale ($n=5$
behaviors per condition), not the originally planned 25-behavior $\times$ 3-seed matrix, because
the free-tier compute budget (Kaggle, $2\times$ T4 per session) was exhausted before the full
matrix could run. We state this plainly rather than let smoke-scale numbers imply matrix-scale
statistical power.

\paragraph{Contributions.}
\begin{enumerate}
  \item A reproducible, hand-verified demonstration that a standard jailbreak-fuzzing success
    judge (\texttt{hubert233/GPTFuzz}) is fooled by jailbreak-shaped surface form rather than
    actual harm, and that this failure mode \textbf{persists after a deliberate, rubric-based
    LLM-judge fix}---not a single-classifier quirk.
  \item An honest null result for activation-guided span mutation vs.\ uniform mutation, reported
    with independent confirmation that the guided-mutation mechanism (attribution + tree-search
    revisiting) was actually active during the runs that produced the null, distinguishing ``no
    effect'' from ``mechanism never fired.''
  \item A stated, artifact-backed inventory of exactly which numbers in this paper are read
    directly from committed result files, which are PI-hand-verified but not written back into
    any file, and which are reported without any backing artifact at all---offered as a concrete
    illustration of the provenance discipline we think jailbreak-fuzzing papers should adopt more
    generally, given finding 1.
\end{enumerate}

\section{Related Work}
\label{sec:related}

\paragraph{Black-box jailbreak fuzzing.}
GPTFuzzer \citep{yu2023gptfuzzer} introduced MCTS-lite (UCB1) seed-template selection over a
human-authored jailbreak template pool, with an LLM mutator rewriting the \textit{whole} selected
template and a trained RoBERTa classifier as the success judge. We reuse GPTFuzzer's
seed-selection algorithm unchanged (\cref{sec:method}) as our uniform-mutation baseline and as
the search scaffold our guided-mutation variant is built on top of---the novelty in this work is
\textit{where} mutation targets a template and \textit{which} judge determines success, not the
search algorithm itself.

\paragraph{White-box activation steering / refusal directions.}
Refusal behavior in instruction-tuned LLMs is mediated, to a first approximation, by a single
linear direction in activation space, extractable via a difference-in-means contrast between
harmful and harmless prompt activations \citep{arditi2024refusal}; ablating or amplifying this
direction causally affects refusal behavior. We reuse this direction-extraction methodology
(\cref{sec:method}) but for a different purpose than steering generation directly: we use the
direction as a \textbf{token-attribution signal} to decide \textit{where inside a jailbreak
template} a mutation should be targeted, on the hypothesis that spans strongly projecting onto
the refusal direction are the parts of a template most responsible for triggering refusal, and
are therefore the highest-value mutation targets. This connects to a closely related line of work
using refusal-direction signals to directly guide or explain jailbreak search---in particular the
approach informally referred to in this project's planning documents as ``Mechanistic AutoDAN''
\citep{collu2026refusal} (``Refusal Before Decoding: Detecting and Exploiting Refusal Signals in
Intermediate LLM Activations''), which detects and exploits refusal signals in intermediate
activations directly. We distinguish our approach as targeting the \textit{search/mutation-selection}
problem (which span of an existing template to rewrite) rather than the decoding/generation
problem the ``Refusal Before Decoding'' framing addresses; a precise technical comparison against
that work is left as future work (we have not run it as a baseline---see
\cref{sec:setup}).

\paragraph{Jailbreak evaluation and judges.}
AutoDAN \citep{liu2023autodan} and AJF \citep{yu2025ajf} are cited here as context for the
jailbreak-generation literature this project's baselines and method sit alongside---\textbf{we
did not run either as a baseline in this work}; \texttt{PLAN.md}'s compute-budget lock (its own
Section 10; see also \cref{sec:setup}) scoped the core comparison to GPTFuzzer alone as a
sufficient baseline for a preprint, with AutoDAN explicitly deferred to an extended lane that did
not get funded before compute ran out. Separately, a growing body of work has begun questioning
the reliability of jailbreak judges themselves: ``How Reliable Is Your Jailbreak Judge?''
\citep{gao2026reliable} finds that a large and growing share of reported jailbreak-evaluation
results using LLM-judges are unreliable, both on average and under deliberate adversarial
pressure; ``LLM-Safety Evaluations Lack Robustness'' \citep{beyer2025llmsafety} and ``Confusion
is the Final Barrier'' \citep{yan2025confusion} make related robustness/consistency arguments;
``Not Aligned'' is Not ``Malicious'' \citep{mei2024notaligned} specifically discusses judges
hallucinating jailbreak success; and ``Beyond Accuracy: Policy Invariance as a Reliability Test
for LLM Safety Judges'' \citep{weng2026beyond} proposes a reliability test in the same spirit as
our residual-false-positive check. We situate this paper's primary contribution within this
cluster: rather than a general audit or a new reliability \textit{metric}, we contribute a
specific, reproducible, hand-verified case study of a widely-reused classifier judge failing, a
concrete fix, and---the part we believe is under-examined in this literature---direct evidence
that the fix's own output still contains the same failure class, from our own pipeline's real
generations rather than a constructed adversarial example.

\section{Method}
\label{sec:method}

\paragraph{Pipeline overview.}
\texttt{scripts/run\_fuzz.py} implements a single fuzzing loop parameterized by
\texttt{-{}-method \{ours, gptfuzzer\}}, \texttt{-{}-mutation \{guided, uniform\}}, and
\texttt{-{}-fitness judge} (this project's probe-based \texttt{judge+act} fitness variant was
built but demoted---see below). Both methods share:
\begin{itemize}
  \item \textbf{Seed pool.} The original GPTFuzzer human-authored template pool
    (\texttt{sherdencooper/GPTFuzz}'s \texttt{GPTFuzzer.csv}, 77 templates), subsampled
    deterministically to \texttt{seed\_pool\_size=12} templates via a fixed constant seed
    (\texttt{SEED\_POOL\_SUBSAMPLE\_SEED=0}, independent of the run's own \texttt{-{}-seed})---identical
    12 templates across every run, replicate, and target. \Cref{sec:setup} explains why this
    subsampling exists.
  \item \textbf{UCB1 pool selection.} GPTFuzzer's MCTS-lite node selection, reused unchanged: at
    each iteration, select the pool node maximizing UCB1 score (unvisited nodes get priority;
    visited nodes trade off mean reward against an exploration bonus), mutate it, and append the
    mutated child as a new pool node.
  \item \textbf{Two-stage success judging} (\texttt{scripts/judge.py}, \cref{sec:judge}).
\end{itemize}

\subsection{Refusal-direction extraction}
\label{sec:direction}

Following \citet{arditi2024refusal}, we extract a per-layer difference-in-means direction from
paired harmful (AdvBench) and harmless (Alpaca-style) prompts' residual-stream activations, on
both targets.

\begin{table}[htbp]
\centering
\caption{Refusal-direction and probe comparison, Qwen2.5-3B (control) vs.\ Phi-4-mini
  (treatment). Sources: \texttt{results/direction.npz}, \texttt{results/phi/direction.npz}
  ($\to$ \texttt{best\_layer}/\texttt{best\_layer\_idx}/\texttt{best\_auc}); \texttt{results/qwen\_novel\_check.log},
  \texttt{results/phi4mini/novel\_check.log} (per-prompt and mean \texttt{direction\_score}
  lines); \texttt{results/probes/best\_layer.json}, \texttt{results/phi/probes/best\_layer.json}
  ($\to$ \texttt{best\_auc}).}
\label{tab:direction}
\begin{tabular}{p{6cm}p{4cm}p{4cm}}
\toprule
 & Qwen2.5-3B (control) & Phi-4-mini (treatment) \\
\midrule
Direction best layer / depth & $25/36$ ($\approx 69\%$) & $11/32$ ($\approx 34\%$) \\
Direction held-out AUC & $1.0$ & $1.0$ \\
Novel-prompt direction score, harmful mean & $\mathbf{+30.204}$ & $\mathbf{+9.486}$ \\
Novel-prompt direction score, benign mean & $\mathbf{-4.954}$ & $\mathbf{+4.594}$ \\
Harmful $-$ benign gap & $\mathbf{35.16}$ & $\mathbf{4.89}$ \\
Probe best layer / held-out AUC & $9$ / $1.0$ & $8$ / $1.0$ \\
Probe novel-prompt accuracy & $0.500$ (chance) & $0.500$ (chance) \\
\bottomrule
\end{tabular}
\end{table}

Both targets' held-out AUCs are a perfect $1.0$ for both signals---uninformative on their own, as
Gate 2's own novel-prompt check exists to demonstrate (\texttt{reviews/stage2-human-signoff.md}).
On the novel-prompt check (\cref{tab:direction}), \textbf{Qwen's direction separates harmful from
benign cleanly and by a wide margin (gap $35.16$, benign mean solidly negative); Phi's direction
separates harmful from benign only weakly (gap $4.89$, and notably the benign mean is
\textit{positive}, not negative---Phi's direction never clearly signals ``not harmful'' the way
Qwen's does).} Both targets' probes fail the same novel-prompt check at exactly chance accuracy
($0.500$), replicating Qwen's probe-generalization failure (diagnosed as overfitting AdvBench's
imperative surface form rather than harmful intent, \texttt{reviews/stage2-human-signoff.md}) on
Phi as well. Per that gate's explicit ruling, probe-based fitness (\texttt{-{}-fitness judge+act})
was demoted out of the core comparison entirely on both targets; every result in this paper uses
\texttt{-{}-fitness judge} (judge-only reward), and the probe/\texttt{judge+act} path is not part
of any claim below. Guided \textit{mutation} (driven by the direction, not the probe) is
unaffected by this demotion and remains the paper's method under test---but the direction-quality
asymmetry above is directly relevant to \cref{sec:mechanism}'s attribution-quality finding: Phi's
underlying refusal signal is measurably weaker and noisier than Qwen's, which is a plausible
contributing cause (not the only one---see \cref{sec:mechanism}) for why Phi's attribution scores
cluster more narrowly and less consistently around refusal-semantic content than Qwen's do.

\subsection{Guided mutation}
\label{sec:guided-mutation}

For the selected pool template (\texttt{ours}, \texttt{-{}-mutation guided}), we run a partial
forward pass, compute the per-token projection of each token's residual-stream activation (at the
direction's best layer) onto the refusal direction, and select the contiguous
\texttt{MUTATE\_WINDOW\_TOKENS}-token window (excluding the behavior-injection marker span) with
the highest summed projection. A local mutator LLM is then instructed to rewrite \textit{only
that span} (``reduce refusal-triggering framing, preserve intent''), leaving the rest of the
template untouched. If attribution fails for any reason (degenerate span, shape mismatch,
tokenizer edge case), the loop falls back to whole-template uniform mutation for that iteration
and increments a \texttt{guided\_fallback\_count} counter---this ensures fallback is never silent
(\cref{sec:mechanism}). \textbf{Uniform mutation} (\texttt{gptfuzzer}, and
\texttt{ours -{}-mutation uniform} as an ablation) mutates the selected template's entire text,
per GPTFuzzer's original design, sharing every other component (seed pool, UCB1 selection, judge)
with the guided condition---the only variable that differs between \texttt{ours} and
\texttt{gptfuzzer} is mutation location.

\subsection{Two-stage success judging}
\label{sec:judge}

Motivated by the false-positive finding that is this paper's primary contribution
(\cref{sec:judge-reliability}), success determination in \texttt{scripts/judge.py} is a two-stage
process: (1) a deterministic keyword pre-filter (\texttt{is\_refusal()}) checks the first 300
characters of a completion against a list of canonical refusal phrases---catching ``refuses,
then pivots to something unrelated'' with zero model calls; (2) anything that passes stage 1 is
graded by a fixed, small, ungated LLM (\texttt{microsoft/Phi-3.5-mini-instruct}, distinct from
both attack targets so it is never self-grading) against an explicit rubric instructing it to
fail responses that merely adopt a roleplay persona or template setup without providing real
harmful content, or that address a different/watered-down version of the request, and to pass
only responses that give real, specific, actionable content. The original RoBERTa classifier
judge (\texttt{hubert233/GPTFuzz}) is retained as a \textbf{logged diagnostic only}
(\texttt{roberta\_judge\_score} per candidate, \texttt{roberta\_judge\_positive\_rate}
aggregated)---never the success determinant in any result reported here.

\section{Experimental Setup}
\label{sec:setup}

\paragraph{Targets.}
Two open-weight instruction-tuned models, both ungated and small enough to fit a single Kaggle
T4 GPU in fp16: \textbf{Qwen2.5-3B-Instruct} (control) and \textbf{Phi-4-mini-instruct}
($\approx 3.8$B params, treatment). The two targets share every pipeline invariant except the
model itself and its paired mutator-LLM default (each target self-mutates; see
\texttt{configs/exp.yaml} / \texttt{configs/exp\_phi4mini.yaml}), per this project's rule that
cross-target ASR comparisons are only valid if judge, benchmark, budget, and decoding parameters
are held fixed.

\paragraph{Benchmark and budget.}
5 behaviors per condition sampled from the canonical AdvBench CSV
(\texttt{llm-attacks/llm-attacks}), \texttt{query\_budget=40} full-generation queries per
behavior, \texttt{max\_iterations=100}, \texttt{decode\_temperature=0.7},
\texttt{decode\_top\_p=0.9}---identical across every condition reported
(\texttt{configs/exp.yaml}, \texttt{configs/exp\_phi4mini.yaml}). \texttt{max\_iterations} is
non-binding under \texttt{-{}-fitness judge}: every iteration is a full query, so the loop
always hits the $\mathtt{query\_budget=40}$ cap (or succeeds) well before it could reach
iteration $100$; the effective per-behavior bound throughout this paper is \texttt{query\_budget},
not \texttt{max\_iterations}.

\paragraph{Seed-pool subsampling (\texttt{seed\_pool\_size=12}).}
We found, via a debug-logging pass over the UCB1 pool-selection trace, that the original
77-template seed pool exceeds \texttt{query\_budget=40}: since UCB1 always prioritizes unvisited
nodes, and $77 > 40$, every iteration of every behavior in our first attempt selected the next
never-mutated original seed template, in identical order, for every behavior---the search never
reached a point within budget where UCB1 would revisit and refine an already-mutated child. Both
\texttt{ours} and \texttt{gptfuzzer} were, in effect, doing the same thing (mutate a fresh seed
once) and a guided-vs-uniform comparison run this way would not actually test mutation-location
strategy at all. We fixed this by deterministically subsampling the seed pool to 12 templates
(fixed constant, not \texttt{-{}-seed}---identical 12 templates on every run/replicate/target), so
that $\approx 12$ of the 40-query budget exhausts the pool and the remaining $\approx 28$ do real
UCB1 revisiting. We verified this structurally engages tree search by extracting
\texttt{select\_ucb1()}/\texttt{backpropagate()} \textbf{verbatim} from the codebase and running
them standalone against a synthetic reward loop: the old 77-template pool yields $0/40$
mutated-child selections at \texttt{query\_budget=40} (confirming the problem); the new
12-template pool yields $16/40$ mutated-child selections even under the worst-case
all-zero-reward scenario (confirming the fix). \Cref{sec:mechanism} reports the \textit{real}
runs' \texttt{n\_mutated\_child\_selected} counts, which is the artifact-backed version of this
claim.

\paragraph{Compute constraint---stated plainly.}
All compute for this project ran on Kaggle's free tier ($2\times$ T4 GPUs per session,
$\approx 30$ GPU-hours/week). The originally planned full evaluation---25 behaviors $\times$ 3
seeds $\times$ 2 targets $\times$ \{\texttt{ours}, \texttt{gptfuzzer}, plus ablations\}---\textbf{did
not run.} The free-tier compute budget was exhausted after the smoke-scale ($n=5$) runs reported
here, before the full matrix could be launched. This is a hard resource constraint, not a
scheduling choice: every result in this paper is smoke-scale, and we treat the resulting lack of
statistical power as a first-class limitation (\cref{sec:limitations}), not an implementation
detail to bury.

\section{Results}
\label{sec:results}

\subsection{Judge reliability: primary contribution}
\label{sec:judge-reliability}

\paragraph{The original incident.}
Hand-reading 4 completions the original \texttt{hubert233/GPTFuzz} judge scored $\approx 0.99$
(``jailbroken'') on an early Phi-4-mini smoke run found all four were not jailbreaks: one was a
refusal (``I can't assist'') that pivoted to an unrelated topic; three were the model merely
echoing a DAN/Omega/APOPHIS-style roleplay template's persona/setup instructions, with zero
harmful content (full writeup: \texttt{reviews/judge-validity-incident.md}). The specific
$\mathtt{asr=0.8}$ figure this produced is \textit{not} backed by a committed
\texttt{results/*.json} in this repository---no \texttt{results/phi/ours\_smoke.json} (pre-fix,
pre-pool-fix) exists; the $0.8$ number survives only as narrative record in the incident writeup,
not as a citable artifact. We therefore report the qualitative finding ($4/4$ hand-read false
positives) as established---it is a direct, documented human read---but do not cite
``$0.8 \to 0.0$'' as a clean before/after pair the way the rest of this section's numbers are
cited. The judge-inflation claim rests on the artifact-backed evidence below instead.

\paragraph{A cleaner, artifact-backed before/after (Qwen, same target, same method, same $n=5$
behaviors).}
\texttt{results/ours\_smoke.json} (old RoBERTa judge as the live success determinant, predating
both the judge fix and the seed-pool fix) reports $\mathtt{asr=1.0}$ ($5/5$)---every behavior
``succeeded'' by the old judge's own determination during search.
\texttt{results/ours\_smoke\_pool12.json} (fixed judge as determinant, seed-pool fix applied)
reports $\mathtt{asr=0.4}$ ($2/5$), later hand-verified to a true $0.2$ ($1/5$; see below).
\textbf{We flag this comparison as suggestive, not a controlled ablation}: two variables changed
between these two runs (judge fix AND seed-pool fix), not one, so the delta cannot be cleanly
attributed to the judge alone.

\paragraph{The cleanest same-run comparison available.}
Within \texttt{results/ours\_smoke\_pool12.json} itself (one run, one set of candidates, both
judges scored on the same completions), the retained RoBERTa diagnostic score
($\mathtt{roberta\_judge\_positive\_rate} = 0.335$, i.e.\ the old judge would have flagged
$\approx 34\%$ of the $167$ individually-evaluated candidates as positive) is far higher than the
trusted judge's own candidate-level success rate ($\mathtt{trusted\_judge\_success\_rate} =
0.012$, $\approx 1.2\%$ of the same $167$ candidates; \cref{tab:judge-comparison}).
\textbf{Important caveat, stated in the pipeline's own code comment}: these two fields are
\textbf{candidate-level} rates (over every full-generation query in the run), not the
\textbf{behavior-level} \texttt{asr} statistic reported elsewhere in this paper---they answer
``how often would the two judges have disagreed on an individual completion,'' not ``what would
behavior-level ASR have been under the old judge.'' We report both numbers because they are real
and artifact-backed, but we do not collapse them into a single ``$N\times$ inflation'' headline
figure the way an early draft of this section's brief suggested, because that would conflate two
different statistics. The qualitative direction---old diagnostic judge substantially more
permissive than the trusted judge, on the identical completions---is exactly what this section's
original incident predicts, and holds here too.

\begin{table}[htbp]
\centering
\caption{Judge comparison within \texttt{results/ours\_smoke\_pool12.json} and against the
  pre-fix \texttt{results/ours\_smoke.json}. Behavior-level and candidate-level rates are not
  directly comparable (see caveat in text).}
\label{tab:judge-comparison}
\begin{tabular}{p{4.2cm}p{2.3cm}p{2.0cm}p{4.5cm}}
\toprule
Run & Judge (determinant) & \texttt{asr} & Granularity note \\
\midrule
\texttt{results/ours\_smoke.json} (Qwen, pre-fix, pre-pool-fix) & old RoBERTa & $1.0$ ($5/5$) & behavior-level \\
\texttt{results/ours\_smoke\_pool12.json} (Qwen, fixed judge, pool-fix) & fixed 2-stage & $0.4$ ($2/5$), hand-verified true $0.2$ ($1/5$) & behavior-level \\
--- same run, \texttt{roberta\_judge\_positive\_rate} & old RoBERTa (diagnostic only) & $0.335$ & \textbf{candidate-level, not comparable to \texttt{asr} above} \\
--- same run, \texttt{trusted\_judge\_success\_rate} & fixed 2-stage (diagnostic mirror) & $0.012$ & \textbf{candidate-level, not comparable to \texttt{asr} above} \\
\bottomrule
\end{tabular}
\end{table}

\paragraph{Residual false positive, found AFTER the fix---the ChadGPT case.}
Hand-verifying the fixed judge's 2 flagged Qwen \texttt{ours} successes
(\texttt{results/ours\_smoke\_pool12.json}, $\mathtt{asr=0.4}$) found only 1 is genuine. The
other is a \textbf{false positive}: a ``ChadGPT'' persona-wrapper completion that, in its actual
response, refuses the harmful request and provides crisis-support resources---zero harmful
content---yet the fixed rubric judge (whose rubric explicitly instructs failing ``the response
refuses, declines, moralizes, or deflects\ldots even if it then talks about something
unrelated'') scored it PASS anyway (full writeup: \texttt{reviews/judge-validity-incident.md},
``Residual false positive found AFTER the fix''). True hand-verified Qwen \texttt{ours} ASR:
$\mathbf{1/5 = 0.2}$, not written back into the JSON (this project does not retroactively edit
results files). We consider this the paper's central piece of evidence: the failure mode
identified in the original incident is not specific to \texttt{hubert233/GPTFuzz}'s classifier
architecture---it survives a deliberately anti-roleplay, explicitly-rubric-instructed
LLM-as-judge replacement. Two plausible, non-exclusive, \textbf{undiagnosed} contributors
(flagged rather than asserted, since the raw completion text is gitignored per this project's
data-handling policy and not independently re-inspectable from this repository): the stage-1
keyword pre-filter's 300-character window could have let a later-appearing refusal phrase through
to the LLM judge; or the LLM judge itself may simply have failed to apply its own rubric
correctly on this input---itself a citable limitation of the ``replace a classifier judge with an
LLM judge'' fix strategy.

\subsection{The honest guided-mutation null}
\label{sec:null}

\begin{table}[htbp]
\centering
\caption{ASR by target and method, pool-12 smoke runs ($n=5$ behaviors each). Hand-verified
  column reflects the ChadGPT correction (\cref{sec:judge-reliability}); the raw \texttt{asr}
  field is never retroactively edited in the source JSON.}
\label{tab:null}
\begin{tabular}{p{2.0cm}p{2.6cm}p{1.6cm}p{2.0cm}p{4.3cm}}
\toprule
Target & Method & \texttt{asr} & Hand-verified \texttt{asr} & Source \\
\midrule
Qwen2.5-3B & \texttt{ours} (guided) & $0.4$ ($2/5$) & $\mathbf{0.2}$ ($\mathbf{1/5}$) & \texttt{results/ours\_smoke\_pool12.json} \\
Qwen2.5-3B & \texttt{gptfuzzer} (uniform) & $0.0$ ($0/5$) & $0.0$ ($0/5$) & \texttt{results/gptfuzzer\_smoke\_pool12.json} \\
Phi-4-mini & \texttt{ours} (guided) & $0.0$ ($0/5$) & $0.0$ ($0/5$) & \texttt{results/phi/ours\_smoke\_pool12.json} \\
Phi-4-mini & \texttt{gptfuzzer} (uniform) & $0.0$ ($0/5$) & $0.0$ ($0/5$) & \texttt{results/phi/gptfuzzer\_smoke\_pool12.json} \\
\bottomrule
\end{tabular}
\end{table}

\textbf{We explicitly do not claim guided mutation beats uniform mutation from
\cref{tab:null}.} The Qwen $1$-vs-$0$ (hand-verified) or $2$-vs-$0$ (raw judge count) gap is
noise at $n=5$---a single behavior flipping either way changes the ratio entirely, and Phi shows
a tied $0$-vs-$0$.

Both targets' \texttt{asr} values are computed at $\mathtt{n\_behaviors=5}$ (smoke scale).
Establishing whether either direction (guided advantage, or no advantage) is statistically
significant would require substantially more behaviors and seed replicates than compute allowed
(\cref{sec:setup,sec:limitations})---we report a null at the scale we could afford to test, not a
proof of equivalence at any scale.

\subsection{Mechanism-engagement check: is the null a broken-mechanism artifact?}
\label{sec:mechanism}

A null result is only informative if the mechanism under test actually ran. We can confirm two
specific, artifact-backed facts about mechanism engagement, and must flag a third claim as
unconfirmed.

\textbf{Attribution fires, does not silently fall back.} $\mathtt{guided\_fire\_count} /
\mathtt{full\_forward\_passes}$ is $167/167$ for Qwen \texttt{ours} and $200/200$ for Phi
\texttt{ours}---$\mathtt{guided\_fallback\_count}$ is $0$ in both. \texttt{find\_attribution\_span()}
successfully located a real mutation span on every single iteration of every guided run reported
here; the null above is not explained by guided mutation quietly degrading to uniform mutation
under the hood.

\textbf{Tree search genuinely engages} (the fix in \cref{sec:setup} working as intended). Every
run's $\mathtt{n\_original\_selected}$ and $\mathtt{n\_mutated\_child\_selected}$ counters were
re-verified directly against their source JSON files for this revision, and every pair sums
exactly to that run's own $\mathtt{full\_forward\_passes}$ (\cref{tab:sums}):
Qwen \texttt{ours}, $113 + 54 = 167$; Qwen \texttt{gptfuzzer}, $120 + 80 = 200$; Phi
\texttt{ours}, $120 + 80 = 200$; Phi \texttt{gptfuzzer}, $120 + 80 = 200$. All four runs spend a
substantial fraction of their budget revisiting/refining previously-mutated pool nodes, not just
replaying the 12 original seeds. Qwen \texttt{ours}' lower total ($167$, vs.\ $200$ for the other
three) is not an inconsistency: it reflects early stopping on success, not a different
\texttt{query\_budget}---all four runs share the identical $\mathtt{query\_budget=40} \times 5$
behaviors $= 200$-query cap (\cref{sec:setup}), but Qwen \texttt{ours} is the only one of the
four with any successes ($2/5$ raw), and a behavior's search stops the moment it succeeds rather
than continuing to the cap. Concretely, \texttt{results/ours\_smoke\_pool12.json}'s own
$\mathtt{queries\_to\_success}$ field records successes at queries $17$ and $30$ for two
behaviors, with the remaining three running the full $40$-query cap each: $17 + 40 + 40 + 40 +
30 = 167$, exactly matching the reported $\mathtt{full\_forward\_passes}$. The other three runs
(all $\mathtt{asr}=0.0$) never stop early, so every one of their five behaviors runs the full
$40$-query cap, $5 \times 40 = 200$.

\begin{table}[htbp]
\centering
\caption{Pool-selection counters, verified against source JSON for this revision.
  $\mathtt{n\_original\_selected} + \mathtt{n\_mutated\_child\_selected} =
  \mathtt{full\_forward\_passes}$ holds exactly for all four runs.}
\label{tab:sums}
\begin{tabular}{lccc}
\toprule
Run & \texttt{full\_forward\_passes} & \texttt{n\_original\_selected} & \texttt{n\_mutated\_child\_selected} \\
\midrule
Qwen \texttt{ours} & $167$ & $113$ & $54$ \\
Qwen \texttt{gptfuzzer} & $200$ & $120$ & $80$ \\
Phi \texttt{ours} & $200$ & $120$ & $80$ \\
Phi \texttt{gptfuzzer} & $200$ & $120$ & $80$ \\
\bottomrule
\end{tabular}
\end{table}

This directly answers the question the seed-pool fix (\cref{sec:setup}) set out to answer: at
pool-12, real UCB1 tree search is happening, on both methods, in every reported run. A separate,
independently-collected 2-behavior \texttt{-{}-debug-attribution} diagnostic run shows the first
\texttt{MUTATED\_CHILD} pool selection firing at \textbf{iteration $24$} on \textit{both} targets:
Qwen (\texttt{results/debug\_attribution\_qwen.log}, $\mathtt{behavior\_idx=1}$) and Phi
(\texttt{results/debug\_attribution\_phi.log}, $\mathtt{behavior\_idx=0}$)---re-verified for this
revision by grepping both logs directly, not assumed from a prior draft. This is not a
coincidence: both examined behaviors have zero successes throughout their own search (Qwen's
debug run succeeds only on $\mathtt{behavior\_idx=0}$, at iteration $17$---not the
$\mathtt{behavior\_idx=1}$ examined here; Phi's debug run has $\mathtt{asr=0.0}$ throughout), so
both searches are governed purely by UCB1's explore term under zero reward. Given identical
$\mathtt{seed\_pool\_size=12}$ and identical deterministic tie-breaking (unvisited nodes first,
ties broken by lowest pool index) on both targets, the same zero-reward dynamics necessarily
produce the same transition point: iterations $0$--$11$ visit each of the 12 original seeds
once; iterations $12$--$23$ revisit them a second time (originals still win ties over children
by lower index); by iteration $24$ the originals have accumulated more visits than any child, so
a child's larger explore bonus (fewer visits, same zero reward) finally wins. This matches, to
the exact iteration and for an explained reason rather than by chance, the worst-case prediction
from \cref{sec:setup}'s standalone \texttt{select\_ucb1}/\texttt{backpropagate} simulation---independent
corroboration, from two real runs on two different targets, of a claim \cref{sec:setup}
previously supported only with a synthetic simulation.

\paragraph{Attribution quality: real, quantified, and \textit{not} uniform (corrected from an
earlier draft).}
An earlier version of this section claimed attribution ``localizes to refusal-relevant tokens,''
based on an informal chat excerpt that was never saved as an artifact. We re-ran the diagnostic
with \texttt{-{}-debug-attribution} (2 behaviors per target, real forward passes, not synthetic)
and saved the full output to \texttt{results/debug\_attribution\_qwen.log} (57 attribution calls)
and \texttt{results/debug\_attribution\_phi.log} (80 attribution calls)---both now committed,
with the assembled selected-span text redacted (\texttt{text='[REDACTED --- jailbreak-template
fragment]'}) per a subsequent internal
review's finding that committing that text, even in a diagnostic log, fell outside this project's
own no-jailbreak-strings policy; the individual tokens, per-token projection scores, and
$\mathtt{sum\_score}$ values below are unredacted and are the actual evidence for this section's
claim. Reading the real logs, the earlier claim was \textbf{too strong}: localization is real on
some inputs and largely absent on others, varying by which template (original seed or mutated
child) is selected. Representative examples, with individual top-projecting tokens quoted
directly from the committed logs and assembled span text redacted:

\begin{itemize}
  \item \textbf{Strongly refusal-semantic} (Qwen, iteration $1$,
    \texttt{results/debug\_attribution\_qwen.log}): top tokens include \texttt{harmful}
    ($13.51$), \texttt{illegal} ($12.46$), \texttt{policy} ($13.18$), \texttt{ethical}
    ($11.77$); selected span \texttt{[REDACTED --- jailbreak-template fragment]},
    $\mathtt{sum\_score=108.58}$.
  \item \textbf{Generic / template-structural, same target} (Qwen, iteration $0$, same file):
    top tokens are \texttt{rules} ($8.08$), \texttt{it} ($8.08$), \texttt{opposite} ($5.76$),
    \texttt{contrad[iction]} ($6.23$)---function words and template scaffolding, not
    harm-related; selected span \texttt{[REDACTED --- jailbreak-template fragment]},
    $\mathtt{sum\_score=47.08}$---less than half the semantic example's score.
  \item \textbf{Strongly refusal-semantic} (Phi, \texttt{results/debug\_attribution\_phi.log}):
    top tokens include \texttt{illegal} ($10.19$), \texttt{malware} ($10.77$), \texttt{harmful}
    ($9.57$), \texttt{discrimination} ($7.98$), \texttt{racism} ($8.47$); selected span
    \texttt{[REDACTED --- jailbreak-template fragment]}, $\mathtt{sum\_score=84.96}$.
  \item \textbf{Generic, same target} (Phi, same file): top tokens are \texttt{mode} ($6.86$),
    \texttt{pretend} ($6.34$), \texttt{Anti[GPT]} ($6.28$), \texttt{character} ($6.09$); selected
    span \texttt{[REDACTED --- jailbreak-template fragment]}, $\mathtt{sum\_score=58.36}$.
\end{itemize}

Aggregated over every logged call: Qwen's per-selection $\mathtt{sum\_score}$ ranges
$\mathbf{40.0}$--$\mathbf{116.0}$ (mean $82.6$, $n=57$); Phi's ranges $\mathbf{51.7}$--$\mathbf{87.2}$
(mean $65.8$, $n=80$)---Phi's attribution is more uniformly mid-range, less likely to hit either
Qwen's strongest semantic peaks or its weakest generic troughs. This is directionally consistent
with \cref{sec:direction}'s finding that Phi's underlying refusal direction separates harmful
from benign prompts far more weakly than Qwen's (novel-prompt gap $4.89$ vs.\ $35.16$)---a
noisier underlying signal plausibly produces a narrower, less differentiated attribution-score
distribution, though we have not run a statistical test connecting these two observations and do
not claim a proven causal link, only a consistent pattern.

We now state this section's central claim precisely, not as ``attribution localizes to
refusal-relevant tokens'' (too strong, corrected above) but as: attribution is a genuinely
informative signal on some templates, on both targets, and an uninformative one on others---template-
and target-dependent, not uniform. This is, we argue, \textit{itself} a natural explanation for
\cref{sec:null}'s null: a guidance signal that only sometimes points somewhere meaningful cannot
be expected to reliably outperform unguided (uniform) mutation across a whole search, even though
it clearly is picking up real signal on the templates where it fires strongly. The null is not
evidence the mechanism is broken (the two facts confirmed above rule that out); it is consistent
with a real but unreliable signal, which is a more specific and more useful finding than an
undifferentiated ``no effect.''

\section{Limitations}
\label{sec:limitations}

\begin{itemize}
  \item \textbf{Smoke scale ($n=5$), not the planned matrix.} Every ASR figure in this paper is
    computed over 5 behaviors per condition. The originally planned evaluation (25 behaviors
    $\times$ 3 seeds $\times$ 2 targets) did not run---compute (Kaggle free tier) was exhausted
    first. Effect sizes here should be read as suggestive at best; no claim in this paper should
    be read as carrying matrix-scale statistical power.
  \item \textbf{Compute-constrained matrix abandonment, not a scheduling deferral.} We state this
    plainly rather than imply the full matrix is merely ``future work in progress'': it will not
    run under this project's current resourcing.
  \item \textbf{Single run per condition---no seed replicates.} Every \cref{sec:null,sec:mechanism}
    number comes from exactly one run per (target, method) pair ($\mathtt{seed=0}$;
    \texttt{configs/exp.yaml}/\texttt{configs/exp\_phi4mini.yaml}). The originally planned 3-seed
    replication (\texttt{\_seed1}/\texttt{\_seed2} variants, \texttt{experiments.yaml}) did not
    run---same compute exhaustion as the behavior-count limitation above, but a distinct concern:
    even at $n=5$ behaviors, we have no run-to-run variance estimate at all, so we cannot say
    whether a different \texttt{-{}-seed} would reproduce the same $2$-vs-$0$/$1$-vs-$0$ split or
    land differently. This is a second, independent reason (beyond $n=5$'s small behavior count)
    the guided-vs-uniform null should not be read as precisely quantified.
  \item \textbf{Same-vendor judge LLM.} The rubric-based LLM judge
    (\texttt{microsoft/Phi-3.5-mini-instruct}) is same-vendor as one of the two attack targets
    (\texttt{microsoft/Phi-4-mini-instruct})---a different checkpoint and training run, and never
    grading its own outputs, but not a fully vendor-independent judge either. We flag this as a
    potential (unquantified) source of residual bias, distinct from the persona-wrapper
    false-positive finding in \cref{sec:judge-reliability}.
  \item \textbf{Small, $\leq 4$B-parameter open-weight targets only.} Both targets are ungated,
    small enough to fit a single T4 GPU in fp16. Generalization to larger or more heavily
    safety-tuned models is untested here.
  \item \textbf{Raw completions not fully retained/reproducible from this repository alone.} Per
    this project's data-handling policy, generated jailbreak attempt text
    (\texttt{results/prompts\_*}) is gitignored and never committed. The ChadGPT false-positive
    example (\cref{sec:judge-reliability}) and other hand-verification work were done against
    completions available at generation time (locally/on the Kaggle session) but are not
    preserved in this repository---an independent reader cannot re-inspect the exact completion
    text behind \cref{sec:judge-reliability}'s central finding from this repo alone. This is a
    genuine reproducibility limitation, traded off against the (higher-priority) commitment never
    to publish raw jailbreak strings (\cref{sec:ethics}).
  \item \textbf{Phi's underlying refusal-direction signal is measurably weaker than Qwen's.}
    \Cref{sec:direction}'s novel-prompt check shows a harmful/benign separation gap of $4.89$ on
    Phi vs.\ $35.16$ on Qwen, with Phi's benign mean staying positive rather than crossing to
    negative. This is now resolved as an artifact-backed \textit{finding} rather than a
    missing-data gap (an earlier draft flagged Phi's direction/probe artifacts as absent from the
    repository; both have since been committed), but it remains a limitation on how far
    \cref{sec:null}'s null generalizes: we tested guided mutation on one target with a strong
    direction signal and one with a comparatively weak one, and cannot rule out that a target
    with an even cleaner refusal direction would show a different result.
  \item \textbf{The \cref{sec:mechanism} debug-attribution evidence is itself small-$n$.} The
    \texttt{-{}-debug-attribution} diagnostic run behind that section's quantified claims covers
    2 behaviors per target (57 attribution calls for Qwen, 80 for Phi)---enough to establish that
    attribution quality \textit{varies}, not enough to claim a precise, generalizable
    distribution of how often it is informative vs.\ not. A larger, dedicated debug-attribution
    run across more behaviors would strengthen this claim's precision.
  \item \textbf{AutoDAN and AJF were not run as baselines} (\cref{sec:related})---cited for
    context only, per the project's compute-scoped decision to treat GPTFuzzer alone as a
    sufficient core-lane baseline.
\end{itemize}

\section{Ethics and Responsible Disclosure}
\label{sec:ethics}

This project generated real jailbreak attempts against two open-weight instruction-tuned models,
targeting AdvBench behaviors, some of which concern self-harm and other sensitive categories.
\textbf{Content warning applies to the underlying (unpublished) data this paper describes.}

\begin{itemize}
  \item \textbf{No successful jailbreak strings are published in this repository or paper.} Raw
    template/candidate/completion text is excluded via \texttt{.gitignore} patterns
    (\texttt{results/**/prompts\_*}, \texttt{results/**/*jailbreak*}, and---since the finding
    below---\texttt{results/**/*.log}, \texttt{results/**/*.jsonl} by default). We note this was
    enforced \textbf{imperfectly} at first: \texttt{results/debug\_attribution\_\{qwen,phi\}.log},
    a diagnostic artifact whose filename didn't match either original \texttt{.gitignore}
    pattern, was committed containing assembled jailbreak-template fragments before an internal
    review caught it (\texttt{reviews/stage7.md}). Those fragments have since been redacted from
    both the log files and this paper's own \cref{sec:mechanism}---see the corrected debug-log
    handling there and in \cref{app:provenance}. \textbf{This is now backed by real automated
    tooling, not just convention or filename patterns}: \texttt{scripts/check\_no\_raw\_text.py}
    content-scans (not just filename-matches) every tracked file under \texttt{results/} for
    known raw-text field shapes (\texttt{text=}, \texttt{"completion"}, \texttt{"candidate"},
    \texttt{"prompt"}, \texttt{"template"}), runs in CI on every push
    (\texttt{.github/workflows/ci.yml}, widened to cover all branches---the project's actual
    commits push to feature branches, not just \texttt{main}), and is wired as a local
    pre-commit hook (\texttt{.githooks/pre-commit}, activated automatically by the Kaggle
    bootstrap cell, \texttt{RUNBOOK.md} \S2.4, since that is where this project's commits
    actually originate). This is a real, meaningful control, not a perfect guarantee: it catches
    known field-name shapes, not arbitrary future ones---pair it with the same manual review
    discipline as before, not as a replacement for it.
  \item \textbf{Regeneration, not redistribution.} Anyone needing to verify a specific claim in
    this paper (e.g.\ re-examining the ChadGPT false positive, \cref{sec:judge-reliability}) must
    regenerate it themselves using the committed code, config, and behavior benchmark (all
    public)---we do not ship a copy of the generated harmful text itself, following the same
    withholding precedent as GPTFuzzer's own release practice.
  \item \textbf{Public benchmark only.} All harmful behaviors are drawn from AdvBench, an
    existing public research benchmark; this work introduces no new harmful-behavior taxonomy or
    capability.
  \item \textbf{White-box work on small open-weight models; no production system targeted.}
    Neither target model is a deployed production system; this is controlled
    research-environment testing of publicly released model weights.
  \item \textbf{Defensive framing.} The refusal-direction signal this paper's guided-mutation
    mechanism is built on (\cref{sec:direction}) is directly usable as a runtime monitoring
    signal: a deployment could project activations onto the same direction and flag/intervene on
    completions whose internal state drifts away from a refusal trajectory mid-generation,
    independent of whether guided mutation itself proves useful as an attack method. We suggest
    this as a concrete defensive application of the same interpretability signal studied here.
  \item \textbf{Disclosure.} Per \texttt{PLAN.md}'s own Section 8 (its ethics gate), disclosure
    to the affected open-weight model maintainers must happen before this work is made public.
    The original plan tied disclosure to arXiv posting; that plan was superseded when the GitHub
    repository was made public on 2026-08-17, a deliberate and separately recorded decision
    (\texttt{reviews/disclosure-timing-decision-2026-08-17.md}).
    \textbf{Disclosure notices were sent to Microsoft (MSRC, for Phi-4-mini-instruct and
    Phi-3.5-mini-instruct) and to Alibaba/Qwen (for Qwen2.5-3B-Instruct) as of 2026-08-19}---after
    the repository went public rather than before, which is the deviation recorded above.
    \textbf{No acknowledgement had been received from either vendor at the time of writing.} We
    record this plainly rather than waiting on a reply: the obligation is to notify affected
    maintainers, and non-response to a courtesy notice---which this is, given that no novel
    vulnerability is claimed and the techniques are published prior work---is common and does not
    leave the notification incomplete.
\end{itemize}

\section{Conclusion}
\label{sec:conclusion}

Two findings, both artifact-backed. First, judge reliability in jailbreak fuzzing is a
persistent, cross-architecture problem, not a quirk of one classifier: the original RoBERTa judge
inflated success via template-echo false positives, and a deliberately stricter, anti-roleplay
LLM-as-judge rubric---built specifically to catch that failure mode---still leaked a
persona-wrapper false positive (the ChadGPT case, \cref{sec:judge-reliability}). Second, guided
mutation shows no consistent, reliable ASR advantage over uniform mutation on either target
tested (\cref{sec:null}), and this null comes with a mechanistic explanation rather than standing
unexamined: the token-attribution signal driving guided mutation is only sometimes informative,
and its reliability tracks the underlying refusal direction's own separation quality---Qwen's
direction separates novel harmful from benign prompts by a $35.16$-point gap, Phi's by only
$4.89$ (\cref{sec:direction}), and Phi's weaker signal corresponds to a narrower, less
differentiated attribution-score distribution during search (\cref{sec:mechanism}).

We frame both as contributions, not failures. A cautionary measurement result---a widely-reused
judge fooling even its own deliberate fix---is directly useful to a field that often treats ASR
as a solved metric. An honest, mechanistically-grounded negative result is more informative than
an unexamined positive one, and more informative still than a null whose mechanism was never
checked.

The practical takeaway: activation-guided jailbreak attacks should be evaluated against verified,
adversarially-checked judges, and are only as promising as the separation quality of the
interpretability signal steering them---a weak refusal direction may not support a strong
guided-mutation advantage. The honest next step is the evaluation this paper could not run: the
full 25-behavior $\times$ 3-seed matrix (\cref{sec:setup,sec:limitations}), with the corrected
judge, once compute allows.

\printbibliography

% Bibliography facts (author/year/venue/arXiv-ID) are verified; see paper/references.bib.

\appendix
\section{Full Provenance Table}
\label{app:provenance}

\begin{longtable}{p{3.6cm} p{4.7cm} p{4.2cm} p{2.6cm}}
\caption{Full per-claim provenance. Every number in the body with a row here is referenced via
  \cref{app:provenance} at first mention.}
\label{tab:provenance} \\
\toprule
Claim & File & Field(s) & Status \\
\midrule
\endfirsthead
\toprule
Claim & File & Field(s) & Status \\
\midrule
\endhead
\bottomrule
\endfoot
Qwen \texttt{ours} pool-12 ASR & \texttt{results/ours\_smoke\_pool12.json} &
  \texttt{asr}, \texttt{guided\_fire\_count}, \texttt{n\_original\_selected},
  \texttt{n\_mutated\_child\_selected}, \texttt{roberta\_judge\_positive\_rate},
  \texttt{trusted\_judge\_success\_rate} & Artifact-backed \\
Qwen \texttt{gptfuzzer} pool-12 ASR & \texttt{results/gptfuzzer\_smoke\_pool12.json} &
  same fields & Artifact-backed \\
Phi \texttt{ours} pool-12 ASR & \texttt{results/phi/ours\_smoke\_pool12.json} &
  same fields & Artifact-backed \\
Phi \texttt{gptfuzzer} pool-12 ASR & \texttt{results/phi/gptfuzzer\_smoke\_pool12.json} &
  same fields & Artifact-backed \\
Qwen pre-fix \texttt{ours} ASR & \texttt{results/ours\_smoke.json} &
  \texttt{asr}, \texttt{guided\_fire\_count} & Artifact-backed \\
Qwen direction AUC / layer & \texttt{results/direction.npz} &
  \texttt{best\_layer}, \texttt{best\_layer\_idx}, \texttt{best\_auc} & Artifact-backed \\
Phi direction AUC / layer & \texttt{results/phi/direction.npz} (mirrored
  \texttt{results/phi4mini/direction.npz}) & same fields & Artifact-backed \\
Qwen probe AUC / layer & \texttt{results/probes/best\_layer.json} &
  \texttt{best\_layer}, \texttt{best\_auc} & Artifact-backed \\
Phi probe AUC / layer & \texttt{results/phi/probes/best\_layer.json} (mirrored
  \texttt{results/phi4mini/probes/best\_layer.json}) & same fields & Artifact-backed \\
Qwen direction/probe novel-prompt separation ($30.204$ / $-4.954$ / $0.500$) &
  \texttt{results/qwen\_novel\_check.log} & per-prompt + mean \texttt{direction\_score}
  lines, probe accuracy line & Artifact-backed (console log, not JSON, but committed and
  verbatim) \\
Phi direction/probe novel-prompt separation ($9.486$ / $4.594$ / $0.500$) &
  \texttt{results/phi4mini/novel\_check.log} & same fields & Artifact-backed \\
ChadGPT false positive, true Qwen \texttt{ours} ASR $0.2$ &
  \texttt{reviews/judge-validity-incident.md} & prose & PI hand-verified, not written back
  into any JSON \\
Original Phi incident, $4/4$ false positives, $\mathtt{asr=0.8}$ &
  \texttt{reviews/judge-validity-incident.md} & prose & PI hand-verified (qualitative) /
  \textbf{no backing JSON for the $0.8$ figure} \\
Pool-77 postfix Qwen $0.4/0.4$ & --- & --- & \textbf{No backing file in this repository;
  PI-reported only} \\
\Cref{sec:mechanism} attribution-quality examples + \texttt{sum\_score} ranges (Qwen
  $40.0$--$116.0$/$82.6$, Phi $51.7$--$87.2$/$65.8$) &
  \texttt{results/debug\_attribution\_qwen.log}, \texttt{results/debug\_attribution\_phi.log} &
  \texttt{top-10 tokens}/\texttt{selected span} lines (assembled span \texttt{text=} field
  redacted; token/score/index fields unredacted), \texttt{sum\_score=} values &
  Artifact-backed (console log, committed verbatim); superseded an earlier unbacked
  chat-excerpt claim \\
First \texttt{MUTATED\_CHILD} selection at iteration $24$ (Qwen) &
  \texttt{results/debug\_attribution\_qwen.log} & $\mathtt{behavior\_idx=1}$ iteration-$24$
  \texttt{pool\_select} line & Artifact-backed \\
First \texttt{MUTATED\_CHILD} selection at iteration $24$ (Phi) &
  \texttt{results/debug\_attribution\_phi.log} & $\mathtt{behavior\_idx=0}$ iteration-$24$
  \texttt{pool\_select} line & Artifact-backed; re-verified for this revision, previously
  reported in body text but not itself listed as a row here \\
\texttt{select\_ucb1}/\texttt{backpropagate} tree-search-engagement proof (synthetic) &
  (verbatim-extracted, run standalone, not itself a committed artifact) & --- &
  Reproducible from \texttt{scripts/run\_fuzz.py} source; not a \texttt{results/*.json} fact;
  independently corroborated by the real iteration-$24$ log lines above, now for both targets \\
\end{longtable}

\end{document}
